\begin{workedexample}[Worked Example: Gain of an aperture antenna]
An aperture of physical area $A$ and efficiency $\eta_a$ has gain
$G = \eta_a\,4\pi A/\lambda^2$. A $1\ \text{m}^2$ dish at $10\ \text{GHz}$
($\lambda = 3\ \text{cm}$) with $\eta_a = 0.6$ gives
\[ G = 0.6\cdot\frac{4\pi(1)}{(0.03)^2} = 8.38\times10^{3} = 39.2\ \text{dBi}. \]
Gain rises as $(D/\lambda)^2$, so larger apertures and higher frequencies focus
the beam more tightly.
\end{workedexample}

\begin{keytakeaways}
\begin{itemize}[leftmargin=1.4em,itemsep=2pt]
\item Directivity is peak focusing; gain $= \eta_a\times$ directivity (efficiency included).
\item Aperture gain $G = \eta_a 4\pi A/\lambda^2$; larger $A/\lambda^2$ means higher gain, narrower beam.
\item Isotropic (dBi) and half-wave dipole (dBd) references differ by $2.15\ \text{dB}$.
\end{itemize}
\end{keytakeaways}

\begin{exercises}
\item Express a gain of 100 in dBi.
\item Doubling frequency at fixed dish area changes gain by how many dB?
\item Convert $12\ \text{dBd}$ to dBi.
\end{exercises}

\furtherreading{Balanis~\cite{balanis} (ch.~2); Kraus~\cite{kraus}.}
